\chapter{Оценка гемодинамического лага по временной динамике фМРТ}\label{ch:delay}

В данной главе рассматривается первая постановка задачи декодирования, учитывающая только временную структуру сигнала. Отклик фМРТ возникает с гемодинамическим лагом~$\Delta t$ относительно предъявления стимула, тогда как ЭЭГ отражает нейронную активность практически без задержки, поэтому оценка~$\Delta t$ необходима для согласования временных шкал модальностей. Для ее получения решается обратная к декодированию задача "--- прогноз фМРТ-отклика по визуальному стимулу с помощью линейной авторегрессионной марковской модели. Приводится вычислительный эксперимент на наборе данных с одновременной регистрацией видеостимулов и фМРТ, демонстрирующий отчетливый минимум ошибки восстановления на локализованной зрительной коре. Результаты главы опубликованы в работе~\cite{dorin2024forecasting}.

\section{Постановка задачи}

%% TODO: обратная задача q: Y -> X_f, F(t) ≈ q[y(t - Δt)];
%% гипотеза о гемодинамическом лаге.
%% Основа: раздел 2.2 магистерской диссертации.

\section{Линейная авторегрессионная модель прогноза фМРТ-отклика}

%% TODO: марковское допущение и приращения; сжатие фМРТ; повоксельная линейная регрессия
%% с L2-регуляризацией и замкнутым решением; авторегрессионное восстановление.
%% Основа: раздел 3.1 магистерской диссертации.

\section{Вычислительный эксперимент}

%% TODO: датасет (30 испытуемых, видеостимулы); локализация зрительной коры;
%% зависимость MSE от Δt, минимум ~5 с; согласование с нейрофизиологическими оценками;
%% анализ распределения весов модели.
%% Основа: раздел 4.1 магистерской диссертации.

\section{Выводы по главе}

%% TODO: оценка Δt переносится в главы 3 и 4 (маски активности, синхронизация модальностей).
